
\subsubsection{Study Design}

This study implements a matched-sample observational design testing two hypotheses: (H-E1) documentation gaps exist at hypothesized severity ($\leq 40$\% achieve compliance at T0+90), and (H-M1) repository commit activity correlates with documentation quality (Spearman $\rho \geq$ 0.30). The methodology enables retrospective temporal analysis combined with activity correlation tests.

\textbf{Repository Sampling:} The study specifies sampling N=120 dataset repositories from HuggingFace Datasets Hub (created 2022-01-01 to 2024-12-31, $\geq 10$ stars), stratified by year to control temporal trends. The current implementation uses synthetic data matching these specifications. Oversampling to N=120 (target N=100 after T0 detection) would mitigate the risk of insufficient successful T0 detections in production deployment.

\textbf{Sampling Criteria:}
\begin{itemize}
\item Platform: HuggingFace Datasets Hub only
\item Creation Window: 2022--2024 (3-year window)
\item Visibility Threshold: $\geq 10$ stars
\item Repository Type: Dataset repositories only
\item Stratification: Equal representation across years (~40 per year)
\end{itemize}

\subsubsection{Temporal Precedence: T0 Detection Protocol}

To measure documentation ''at initial release,'' the study requires precise T0 (release timestamp) for each repository. The study implements a 3-tier fallback strategy:

\textbf{Tier 1 (Preferred):} Release tag timestamp via GitHub API  
\textbf{Tier 2 (Fallback):} First dataset commit matching upload pattern  
\textbf{Tier 3 (Last Resort):} Repository creation date

The protocol queries each tier sequentially, accepting the first successful match. This design balances precision (prefer tier 1) with coverage (guarantee tier 3).

\subsubsection{Documentation Completeness Measurement}

The study assesses documentation quality using DCS\_3 (Documentation Completeness Score, 3-component subset), adapted from Rondina et al. Each component is scored on a 0-1 scale:

\textbf{Component 1: Data Collection Context} (0-1)  
\begin{itemize}
\item 1.0: README documents data sources AND collection methodology
\item 0.5: Partial (sources OR methodology, not both)
\item 0.0: No mention of sources or methodology
\end{itemize}

\textbf{Component 2: Preprocessing Transparency} (0-1)  
\begin{itemize}
\item 1.0: Documents ALL of: cleaning, augmentation, split ratios
\item 0.5: Documents 1-2 aspects
\item 0.0: No preprocessing documentation
\end{itemize}

\textbf{Component 3: Licensing Clarity} (0-1)  
\begin{itemize}
\item 1.0: Clear LICENSE file present
\item 0.5: License mentioned in README but no LICENSE file
\item 0.0: No license information
\end{itemize}

\textbf{Total DCS\_3:} Sum of three components (range: 0-3). Compliance threshold: DCS\_3 $\geq$ 2.4 (80\% of maximum score).

For each repository, the protocol specifies cloning the state at T0+90 days, ensuring measurement reflects documentation present 90 days post-release. The current implementation uses synthetic data approximating this measurement.

\textbf{Inter-Rater Reliability:} The protocol specifies that 20\% of repositories (N=20) should be dual-coded by independent coders, with Cohen's kappa ($\kappa$) computed for binary compliance. Quality gate: $\kappa \geq$ 0.70. The synthetic data implementation yields $\kappa$ = 1.00 (perfect agreement) to validate the gate logic.

\subsubsection{Activity Metric Collection}

To test the mechanism hypothesis (H-M1), the study specifies collecting three activity dimensions via GitHub API for each repository, measured over the T0 to T0+90 window:

\textbf{Metric 1: Commits per Month}  
Count total commits in T0--T90 window, divide by 3 (90 days = 3 months).

\textbf{Metric 2: Unique Contributors}  
Extract unique commit author logins in T0--T90 window.

\textbf{Metric 3: Median Issue Response Time}  
For issues created in T0--T90 window, calculate time between creation and first response, compute median.

Additionally, repository age at T0+90 is recorded for use as a control variable in partial correlation. The current implementation generates synthetic values matching expected distributions.

\subsubsection{Statistical Analysis}

\textbf{Existence Hypothesis (H-E1)}

\textbf{Primary Test:} Binomial proportion test with Wilson score confidence interval.  
\begin{itemize}
\item Null hypothesis (H0): $\pi \geq$ 0.70 ($\geq 70$\% achieve DCS\_3 $\geq$ 2.4)
\item Alternative hypothesis (H1): $\pi$ < 0.70
\item Gate criterion: 95\% CI upper bound < 0.60
\end{itemize}

\textbf{Secondary Test:} Chi-square goodness-of-fit for component breakdown.

\textbf{Mechanism Hypothesis (H-M1)}

\textbf{Primary Test:} Spearman rank correlation between commits/month and DCS\_3.  
\begin{itemize}
\item Null hypothesis (H0): $\rho \leq$ 0.10 or p $\geq$ 0.05
\item Alternative hypothesis (H1): $\rho \geq$ 0.30 and p < 0.05 (one-tailed)
\item Gate criterion: $\rho \geq$ 0.30 with p < 0.05
\end{itemize}

\textbf{Secondary Test:} Partial correlation controlling for repository age.  
\begin{itemize}
\item Gate criterion: partial $\rho \geq$ 0.25 and p < 0.05
\end{itemize}

\textbf{Mechanism Specificity Tests:} Spearman correlation for contributors and issue response metrics to test whether correlation is commit-specific or generalizes to all activity dimensions.

\subsubsection{Implementation}

The study was implemented as a proof-of-concept validation using synthetic data matching expected distributions. The synthetic data generator was designed to produce: (1) approximately 35-40\% initial compliance rates, (2) non-uniform component scores with licensing weakest, and (3) strong correlation between commits and documentation quality. Production deployment requires replacing synthetic data generation with actual HuggingFace Hub API and GitHub API calls, following the protocols specified above. Manual DCS\_3 coding would be performed by trained coders using a standardized template.

All analysis code uses standard statistical libraries (scipy.stats, statsmodels) and follows established protocols for binomial proportion tests, chi-square tests, and Spearman correlation with bootstrap confidence intervals.
